% from ega4-9-fr.tex
\oldpage[IV-3]{91}

donc on peut remplacer $S$ par $S'$ et $X$ par $X'$, autrement dit supposer que
$S$ est \emph{intègre} et que $\eta=f(x)$ est son point générique.

Or, si l'on pose $Z=\operatorname{Supp}(\mathcal F)$, le même raisonnement que
dans la démonstration de (9.9.1) montre que, si $Z_i$ sont les composantes
irréductibles de $Z$ contenant $x$, on peut supposer que ce sont exactement les
composantes irréductibles de $Z$ rencontrant $\overline{\{x\}}$ et que
$(Z_i)_s\cap\overline{\{x\}}\neq\varnothing$ pour tout $s\in S$. La conclusion
résulte alors de (9.8.3) et (9.8.6).

(iii) On se ramène comme dans (ii) au cas où $S$ est noethérien et intègre et
$\eta=f(x)$ son point générique, en utilisant (4.2.7). Comme dans la
démonstration de (9.9.1), on est ramené à montrer qu'il existe un voisinage $V$
de $x$ dans $\overline{\{x\}}$ tel que, pour tout $x'\in V$, l'ensemble des
dimensions des cycles premiers associés à $\mathcal F_{f(x')}$ qui contiennent
$x'$ soit \emph{le même}. Or, si $Z'_i$ sont les adhérences dans $X$ des cycles
premiers associés à $\mathcal F_\eta$ qui contiennent $x$ (cf. (9.8.2)), il
résulte de (9.8.3) et (9.5.1) que pour tout $s$ assez voisin de $\eta$, les
cycles premiers associés à $\mathcal F_s$ qui rencontrent
$\overline{\{x\}}$ sont des composantes irréductibles des $(Z'_i)_s$ et, en
vertu de (9.5.6), toutes ces composantes irréductibles ont même dimension égale
à $\dim((Z'_i)_\eta)$, d'où la conclusion.

(iv) On raisonne comme dans (iii), en utilisant cette fois (4.7.11). Avec les
mêmes notations que dans (iii), les cycles premiers associés à $\mathcal F_s$
qui sont des composantes irréductibles de $(Z'_i)_s$ sont immergés si et
seulement si $(Z'_i)_\eta$ est un cycle premier associé immergé de
$\mathcal F_\eta$. On conclut déjà que si $x$ appartient à un (resp.
n'appartient à aucun) cycle premier associé immergé de $\mathcal F_\eta$,
l'ensemble des $x'\in\overline{\{x\}}$ tels que $x'$ appartienne à un (resp.
n'appartienne à aucun) cycle premier associé immergé de $\mathcal F_{f(x')}$
est un voisinage de $x$ dans $\overline{\{x\}}$. La conclusion résulte de cette
remarque, de la caractérisation des points où un Module est géométriquement
réduit (4.7.10), et de (ii).

(v) Compte tenu de (4.7.11), on se ramène encore au cas où $S$ est noethérien
et intègre et où $\eta=f(x)$ est son point générique. Soit $Y$ un
sous-préschéma fermé de $X$ ayant $\operatorname{Supp}(\mathcal F)$ pour espace
sous-jacent. Raisonnant comme au début de la démonstration de (9.7.7), on voit
qu'il existe une $A$-algèbre intègre finie $A'$ (de sorte que si
$S'=\operatorname{Spec}(A')$, le morphisme $S'\to S$ est fini et surjectif)
telle que si $Y'=Y_{(S')}$ et si $\eta'$ est le point générique de $S'$, les
composantes irréductibles de $Y'_{\eta'}$ soient géométriquement irréductibles.
D'autre part, si $X'=X_{(S')}$, le morphisme projection $h:X'\to X$ est fini
et surjectif, donc fermé~; par suite, si $x'$ est un point de $X'$ au-dessus de
$x$, on a $h(\overline{\{x'\}})=\overline{\{x\}}$ et l'image par $h$ d'un
voisinage de $x'$ dans $\overline{\{x'\}}$ est un voisinage de $x$ dans
$\overline{\{x\}}$. Compte tenu de (4.7.11) et posant
$\mathcal F'=\mathcal F_{(S')}$, on est donc ramené à prouver que si
$\mathcal F'_{\eta'}$ est (resp. n'est pas) géométriquement intègre au point
$x'$, l'ensemble des $y'\in\overline{\{x'\}}$ tels que
$\mathcal F'_{f'(y')}$ (où $f'=f_{(S')}:X'\to S'$) soit (resp. ne soit pas)
géométriquement intègre au point $y'$ est un voisinage de $x'$ dans
$\overline{\{x'\}}$. On peut donc se borner au cas où $X'=X$ et où les
composantes irréductibles de $Y_\eta$ sont géométriquement irréductibles~; si
l'on désigne par $Z_i$ des parties fermées de $X$ telles que les
$Z_i\cap X_\eta$ soient les

